\documentclass[10pt]{article}
\usepackage[utf8]{inputenc}
\usepackage[T1]{fontenc}
\usepackage{amsmath}
\usepackage{amsfonts}
\usepackage{amssymb}
\usepackage[version=4]{mhchem}
\usepackage{stmaryrd}
\usepackage{bbold}
\usepackage{graphicx}
\usepackage[export]{adjustbox}
\graphicspath{ {./images/} }

\begin{document}
Second Welfare Theorem (SWT)

\section*{Preliminaries SWT}
\begin{itemize}
  \item Let $B \subset \mathbb{R}, a \in \mathbb{R}$, then
  \item $\sup B \leq a \Longrightarrow[(\forall b \in B) b \leq a]$
  \item $\inf B \geq a \Longrightarrow[(\forall b \in B) b \geq a]$
  \item Let $A_{1}, A_{2}$ be subsets of a vector space $X$. Then,
  \item $A_{1}+A_{2}=\left\{x_{1}+x_{2}: x_{1} \in A_{1}, x_{2} \in A_{2}\right\}$.
  \item $\operatorname{co} A_{1}$ is the convex hull of $A_{1}$
  \item If $A_{1}, A_{2}$ are convex, then $A_{1}+A_{2}$ is convex.
  \item $p \cdot A_{1}=\left\{p \cdot x_{1}: x_{1} \in A_{1}\right\}$.
\end{itemize}

\section*{Separating Hyperplane Theorem}
Theorem 1 (Minkowski) If $B, C \subseteq \mathbb{R}^{L}$ are convex and $B \cap C=\emptyset$, then $\exists p \in \mathbb{R}^{L}, p \neq 0$, such that $\sup p \cdot B \leq \inf p \cdot C$.

In particular, $b \in B, c \in C \Longrightarrow p \cdot b \leq p \cdot c$.

Second Welfare Theorem

Second Welfare Theorem

Second Welfare Theorem

Second Welfare Theorem

\section*{Weak second welfare theorem}
Theorem 2 (WSWT) Let $\mathcal{E}$ be an Arrow-Debreu economy with convex production sets and complete, transitive, convex, LNS preferences. Let ( $x^{*}, y^{*}$ ) be a Pareto optimal allocation. Then ( $\exists p^{*} \in \mathbb{R}^{L}, p^{*} \neq 0$ ) such that

\begin{enumerate}
  \item $(\forall j)\left(\forall y_{j} \in Y_{j}\right) p^{*} \cdot y_{j}^{*} \geq p^{*} \cdot y_{j}$ (price-taking profit maximization)
  \item $(\forall i)\left[x_{i} \succ_{i} x_{i}^{*} \Longrightarrow p^{*} \cdot x_{i} \geq p^{*} \cdot x_{i}^{*}\right]$ (price-taking "quasi preference maximization")
  \item $\sum_{i=1}^{I} x_{i}^{*}=\bar{\omega}+\sum_{j=1}^{J} y_{j}^{*}$
\end{enumerate}

Any $\left(x^{*}, y^{*}, p^{*}\right)$ satisfying 1-3 is a quasi-equilibrium.

\section*{Proof}
Since $\left(x^{*}, y^{*}\right)$ is an allocation, (3) holds. We will prove (1) and (2) in several steps. Let

$$
\begin{aligned}
V_{i} & =\left\{x_{i} \in X_{i}: x_{i} \succ_{i} x_{i}^{*}\right\} \\
V & =\sum_{i=1}^{I} V_{i} \\
Y & =\sum_{j=1}^{J} Y_{j}
\end{aligned}
$$

Claim (i). $V, Y$, and $\left\{V_{i}\right\}_{i=1}^{I}$ are convex. It follows from convexity of preferences and production sets and because the sum of convex sets is a convex set.

Claim ii. $V \cap(\bar{\omega}+Y)=\emptyset$

\begin{itemize}
  \item Suppose not. Then $\exists \sum_{i=1}^{I} x_{i} \in V: \sum_{i=1}^{I} x_{i} \in(\bar{\omega}+Y)$.
  \item Hence $\left(x^{*}, y^{*}\right)$ is not PO because for every $i, x_{i} \succ_{i} x_{i}^{*}$, and there is $y \in Y$ such that $(x, y)$ is a feasible allocation.
\end{itemize}

This contradiction proves the claim.

Claim iii. $\exists p^{*} \in \mathbb{R}^{L}, p^{*} \neq 0: \inf p^{*} \cdot \bar{V} \geq \sup p^{*} \cdot(\bar{\omega}+Y)$

\begin{itemize}
  \item $V$ and $\bar{\omega}+Y$ are nonempty, convex and do not intersect. By SHT,\\
$\exists p^{*} \in \mathbb{R}^{L}, p^{*} \neq 0: \inf p^{*} \cdot V \geq \sup p^{*} \cdot(\bar{\omega}+Y)$. Thus
\end{itemize}

$$
\begin{aligned}
\left(\sum_{i=1}^{I} x_{i} \in V\right) \wedge & \left(\sum_{j=1}^{J} y_{j} \in Y\right) \Longrightarrow \\
& {\left[p^{*} \cdot \sum_{i=1}^{I} x_{i} \geq p^{*} \cdot\left(\bar{\omega}+\sum_{j=1}^{J} y_{j}\right)\right] }
\end{aligned}
$$

\begin{itemize}
  \item If $z \in \bar{V}$, then $\exists z^{n} \in V: z^{n} \rightarrow z$ (definition of closure)
  \item $p^{*} \cdot z^{n} \geq \sup p^{*} \cdot(\bar{\omega}+Y)$
  \item Continuity of dot product implies $p^{*} \cdot z \geq \sup p^{*} \cdot(\bar{\omega}+Y)$.
\end{itemize}

Recall for future use that $\bar{\omega}+Y \ni \bar{\omega}+\sum_{j} y_{j}^{*}=\sum_{i} x_{i}^{*}$.

Claim iv. $\forall z \in \bar{V}, \forall y \in Y$,\\
$p^{*} \cdot z \geq p^{*} \cdot \sum_{i=1}^{I} x_{i}^{*}=p^{*} \cdot\left(\bar{\omega}+\sum_{j=1}^{J} y_{j}^{*}\right) \geq p^{*} \cdot\left(\bar{\omega}+\sum_{j=1}^{J} y_{j}\right)$\\
Proof. By LNS for $i=1, \ldots, I$,

$$
\forall \epsilon>0, \exists \tilde{x}_{i}^{\epsilon},\left\|\tilde{x}_{i}^{\epsilon}-x_{i}^{*}\right\|<\epsilon \wedge \tilde{x}_{i}^{\epsilon} \in V_{i}
$$

Thus, $\sum_{i=1}^{I} \tilde{x}_{i}^{\epsilon} \in V$

\begin{itemize}
  \item As $\epsilon \rightarrow 0, \forall i, \tilde{x}_{i}^{\epsilon} \rightarrow x_{i}^{*}$ and therefore $\sum_{i=1}^{I} \tilde{x}_{i}^{\epsilon} \rightarrow \sum_{i=1}^{I} x_{i}^{*}$. Thus, $\sum_{i=1}^{I} x_{i}^{*} \in \bar{V}$.
  \item $\sum_{i=1}^{I} x_{i}^{*}=\bar{\omega}+\sum_{j=1}^{J} y_{j}^{*} \in \bar{\omega}+Y$.
  \item Therefore, $\sum_{i=1}^{I} x_{i}^{*} \in \bar{V} \cap(\bar{\omega}+Y)$ and the claim follows.\\
\includegraphics[max width=\textwidth, alt={}, center]{f0229c61-9af6-488f-a6cb-d1d2024f4c46-15_828_760_192_152}\\
\includegraphics[max width=\textwidth, alt={}, center]{f0229c61-9af6-488f-a6cb-d1d2024f4c46-15_832_772_190_967}\\
\includegraphics[max width=\textwidth, alt={}, center]{f0229c61-9af6-488f-a6cb-d1d2024f4c46-15_837_773_190_1790}
\end{itemize}

Claim v. (1) holds, that is to say $(\forall j, j=1, \ldots, J)$,

$$
\left[p^{*} \cdot y_{j}^{*} \geq p^{*} \cdot y_{j}, \forall y_{j} \in Y_{j}\right] \text { (profit maximization) }
$$

Proof. Fix a firm $j$.

\begin{itemize}
  \item $\forall y_{j} \in Y_{j}, y_{j}+\sum_{h=1, h \neq j}^{J} y_{h}^{*} \in Y$.\\
-By claim iv,
\end{itemize}

$$
p^{*} \cdot\left(\bar{\omega}+y_{j}^{*}+\sum_{h=1, h \neq j}^{J} y_{h}^{*}\right) \geq p^{*} \cdot\left(\bar{\omega}+y_{j}+\sum_{h=1, h \neq j}^{J} y_{h}^{*}\right)
$$

\begin{itemize}
  \item Substract the first and third terms to both sides to establish the claim.
\end{itemize}

Claim vi. (2) holds, that is to say

$$
\text { for } i=1, \ldots, I,\left[x_{i} \succ_{i} x_{i}^{*} \Longrightarrow p^{*} \cdot x_{i} \geq p^{*} \cdot x_{i}^{*}\right]
$$

Proof. Fix a consumer $i$.

\begin{itemize}
  \item By claim iv, $p^{*} \cdot\left(x_{i}+\sum_{h=i, h \neq i}^{I} x_{h}^{*}\right) \geq p^{*} \cdot\left(x_{i}^{*}+\sum_{h=i, h \neq i}^{I} x_{h}^{*}\right)$
  \item Substracting the second term to both sides, the claim and the theorem are established.\\
QED
\end{itemize}

Arrow's exceptional economy

\section*{A corollary to the WSWT}
Corollary $\mathbf{1}$ (WSWC) Let $\mathcal{E}$ be an Arrow-Debreu economy that satisfies the hypothesis of Theorem 2 and $p^{*}$ be the price identified in that theorem. If, in addition, preferences are continuous, consumption sets are convex and for $i=1, \ldots, I$ there is a bundle $x_{i}^{\prime} \in X_{i}$ such that $p^{*} \cdot x_{i}^{\prime}<p^{*} \cdot x_{i}^{*}$, then there exists an income transfer $t \in \mathbb{R}^{I}, \sum_{i=1}^{I} t_{i}=0$ such that $\left(x^{*}, y^{*}, p^{*}, t\right)$ is a WE with transfers of $\mathcal{E}$.

\begin{itemize}
  \item Unsatisfactory version of SWT; it uses assumptions on endogenous elements of the model.
\end{itemize}

\section*{Proof [sketch.]}
Theorem 2 (1) and (3) establish that firms maximize profit at the proposed production plan and that aggregate demand equals aggregate supply.

It remains to show preference maximization, that $\forall i$

$$
\forall i, x_{i}^{*} \in D_{i}\left(p^{*}, y^{*}, t_{i}\right) .
$$

Theorem 2 (2) states that

$$
\forall i, x_{i} \succ_{i} x_{i}^{*} \Longrightarrow p^{*} \cdot x_{i} \geq p^{*} \cdot x_{i}^{*}
$$

Claim. $\forall i, x_{i} \succ_{i} x_{i}^{*} \Longrightarrow p^{*} \cdot x_{i}>p^{*} \cdot x_{i}^{*}$.

\begin{itemize}
  \item Arguing by contradiction, suppose $\exists x_{i} \succ_{i} x_{i}^{*} \wedge p^{*} \cdot x_{i}=p^{*} \cdot x_{i}^{*}$ (using Theorem 2 (2))
  \item Pick $x_{i}^{\prime}$ as in the hypothesis of corollary ( $x_{i}^{\prime}$ costs less than $x_{i}^{*}$ ).
  \item By continuity of preferences
\end{itemize}

$$
\exists \bar{\alpha} \in(0,1): 1 \geq \alpha \geq \bar{\alpha} \Longrightarrow \alpha x_{i}+(1-\alpha) x_{i}^{\prime} \succ_{i} x_{i}^{*}
$$

\begin{itemize}
  \item Since $p^{*} \cdot x_{i}^{\prime}<p^{*} \cdot x_{i}^{*}=p^{*} \cdot x_{i}$, then
\end{itemize}

$$
\forall \alpha \in[0,1), p^{*} \cdot\left[\alpha x_{i}+(1-\alpha) x_{i}^{\prime}\right]<p^{*} \cdot x_{i}^{*}
$$

This is a contradiction of Theorem 2 (2).

It remains to identify transfers $\left\{t_{i}\right\}_{i=1}^{I}, \sum_{i=1}^{I} t_{i}=0$, so that $x_{i}^{*} \in D_{i}\left(p^{*}, y^{*}, t_{i}\right)$.

For $i=1, \ldots, I$, define,

$$
t_{i}=p^{*} \cdot x_{i}^{*}-p^{*} \cdot \omega_{i}-\sum_{j=1}^{J} \theta_{i j} p^{*} y_{j}^{*}
$$

Then

\begin{itemize}
  \item $\sum_{i=1}^{I} t_{i}=0$, so $t=\left(t_{1}, \ldots, t_{I}\right)$ is an income transfer.
  \item By construction, $x_{i}^{*} \in B_{i}\left(p^{*}, y^{*}, t_{i}\right)$. By previous argument any $x_{i}, x_{i} \succ_{i} x_{i}^{*}$ must cost more. Then $x_{i}^{*} \in D_{i}\left(p^{*}, y^{*}, t_{i}\right)$.
\end{itemize}

QED.

\section*{A version of the Second Welfare Theorem}
Theorem 3 If $x^{*}$ is a Pareto optimal allocation in a pure exchange economy, with strongly monotone, continuous, convex preferences, consumption sets are $\mathbb{R}_{+}^{L}$, and $\bar{\omega} \gg 0$, then there exists a price vector $p^{*}$ and an income transfer $t$ such that $\left(x^{*}, p^{*}\right)$ is a Walrasian Equilibrium with transfers provided that $(\forall i) x_{i}^{*} \neq 0$.

\section*{Failures of the SWT}
\begin{itemize}
  \item There are three general types of failures.
  \item Failure for insufficient income.
  \item Failure of LNS.
  \item Failure of convexity, no appropriate prices exist.
\end{itemize}

\section*{Failure of LNS}
Failure of convexity

\section*{Failure of convexity in production}
\begin{center}
\includegraphics[max width=\textwidth, alt={}]{f0229c61-9af6-488f-a6cb-d1d2024f4c46-28_1414_2361_348_147}
\end{center}

\section*{Failure of convexity in production}
\begin{center}
\includegraphics[max width=\textwidth, alt={}]{f0229c61-9af6-488f-a6cb-d1d2024f4c46-29_1414_2358_348_150}
\end{center}

\section*{Interpretation and limitations}
\begin{itemize}
  \item Planner must ensure that the new supporting prices are taken as given. This fails with market power; outcome of market economy is not a WE.
  \item Planner's informational requirements: Identify the PO allocation to be implemented and compute the supporting prices in case transfers are used.
  \item In both points above, planner must know the statistical distribution of preferences and endowments.
  \item To implement transfers, the planner must observe the agents' private information, endowment and preferences. This is difficult; it often prevents lump sum transfers.
  \item When lump-sum transfers are not possible, wealth redistribution is distortionary. Thus, the end result is not the desired PO allocation. This leads to theory of second best.
  \item Transfers must be enforced.
\end{itemize}

Mas-Colell, Andreu, Michael Whinston, and Jerry Green. 1995. Microeconomic Theory. New York, Oxford: Oxford University Press.


\end{document}